\chapter[Open Problems and Research Programme]{Open Problems and Research Programme}\label{ch:sec:open}

The deterministic realization problem is completely characterized in this book.  Several natural extensions remain open.  This chapter organises them into four categories: structural problems concerning the WIS representation itself, geometric problems arising from the posterior manifold, algorithmic and computational problems, and applications to concrete security domains.

\section*{Structural Problems}

Problems concerning the WIS representation, its factorisation theory, and the classification of encoders.

\section{Probabilistic Encoders and Randomised WIS}\label{ch:sec:probabilistic-encoders-and-randomised-wis}

The entire book has assumed a deterministic encoder
$E:\mathcal M\times\mathcal K\to\mathcal C$.  Allowing the encoder
to depend on an additional random seed---or, equivalently, replacing
the deterministic map by a row-stochastic matrix---relaxes the
integrality condition that is the main obstruction to closed-form
construction.

\subsection{The Randomised WIS}\label{ch:sec:the-randomised-wis}

Let $\mathcal S$ be a finite set of seeds, and let
$p(s)$ be a probability distribution on $\mathcal S$.
A randomised encoder is a family of deterministic encoders
$\{E_s\}_{s\in\mathcal S}$ indexed by the seed, applied with
probability $p(s)$.  The resulting multiplicity matrix is the
convex combination
\[
n_{mc} = \sum_{s\in\mathcal S} p(s)\,n_{mc}^{(s)},
\]
where $n_{mc}^{(s)}$ are the multiplicities of the deterministic
encoder $E_s$.  Since each $n_{mc}^{(s)}$ is integral, the
randomised multiplicity $n_{mc}$ is a rational number with
denominator dividing the least common multiple of the seed
probabilities.

The universal optimality condition
$|\mathrm{Cloak}(c,\cdot)| = |\mathcal K|\cdot\Pr[M=m\mid C=c]$
depends only on the averaged multiplicities, not on the individual
encoders.  Therefore a randomised encoder is universally optimal
if and only if the averaged WIS $n_{mc}$ admits a
$(\mathbb Q_{>0},\mathbb Q_{>0})$ factorisation---i.e., the
factorisation exists with rational weights.

\begin{theorem}[Randomised Universal Optimality]
\label{ch:thm:randomised-uod}
Let $\{E_s\}_{s\in\mathcal S}$ be a finite family of deterministic
encoders with multiplicities $n_{mc}^{(s)}$, and let
$p(s)>0$ be rational probabilities.  The randomised encoder
obtained by sampling $s\sim p$ and applying $E_s$ is universally
optimal if and only if the averaged multiplicities
\[
\overline n_{mc} = \sum_s p(s)\,n_{mc}^{(s)}
\]
admit a $(\mathbb Q_{>0},\mathbb Q_{>0})$ WIS factorisation.
In particular, every universally optimal rational WIS can be
realised by a randomised encoder with finite seed set.
\end{theorem}

\begin{proof}
The averaged multiplicity is rational by construction.  The
universal optimality condition is linear in the multiplicities,
so the condition reduces to factorisability of $\overline n_{mc}$.
Conversely, given a rational factorisation, the Apollonian
construction~\cite{biondi2018apollonian} produces a deterministic
encoder for each integer scaling; the randomised encoder is
obtained by mixing these with appropriate probabilities.
\end{proof}

\begin{corollary}[Rational Approximation]\label{ch:cor:rational-approx}
Every real WIS weight that is not realisable by a deterministic
encoder can be approximated arbitrarily well by a randomised
encoder whose seed set size grows with the desired precision.
The approximation error (measured by the Universal Optimality
Defect) decays as $O(1/|\mathcal S|)$.
\end{corollary}

\begin{proof}
Approximate the real weights by rational numbers and apply
Theorem~\ref{ch:thm:randomised-uod}.  The $O(1/|\mathcal S|)$
rate follows from standard rational approximation bounds.
\end{proof}

\paragraph{Open problems.}
\begin{itemize}
\item \textbf{Optimal seed distribution.}  Given a target WIS
weight, what is the minimal seed set size $|\mathcal S|$
required to realise it within a given approximation tolerance?
\item \textbf{Integrality gap.}  Given rational weights, what is
the minimal integer scaling factor $T$ such that $T\cdot w$ and
$T\cdot\beta$ are integral?
\end{itemize}

\section{Continuous Weighted Incidence Structures}\label{ch:sec:open-continuous-wis}

Extend the WIS framework from finite to continuous message and
ciphertext spaces.  When the incidence graph is replaced by a
continuous bipartite structure (e.g. a kernel), the balance equation
$B\beta = Kw$ becomes an integral equation.

\paragraph{Open problems.}
\begin{itemize}
\item \textbf{Infinite message spaces.}  Extending the WIS framework to
countable or continuous message spaces requires measure-theoretic
formulations of the quotient geometry and the UOD.
\item \textbf{Continuous WIS.}  When the incidence is continuous, the
balance equation becomes an integral equation.  Existence of solutions
and their interpretation as encoding strategies is unexplored.
\item \textbf{Approximation theory.}  When induced multiplicity ratios
are irrational, no exact finite encoder exists.  Rational approximation
and the trade-off between quality and key length merit investigation.
\end{itemize}

\section{Random WIS and Average-Case Analysis}\label{ch:sec:open-random-wis}

What does a typical WIS look like?  For random incidence relations,
the UOD is typically large.  Understanding its distribution provides
a baseline for evaluating real-world encryption systems.

\paragraph{Open problems.}
\begin{itemize}
\item \textbf{Typical UOD.}  Expected value of $\mathcal D$ for a
random WIS with fixed $|\mathcal M|,|\mathcal C|$?
\item \textbf{Phase transitions.}  Density threshold below which
universal optimality is impossible with high probability?
\item \textbf{Random graph models.}  Do the spectral properties of the
posterior Laplacian for random WIS match those of random graphs?
\end{itemize}

\section{Categorical Formulations}\label{ch:sec:open-categorical}

The WIS construction, the gauge quotient, and the Hilbert-space
structure admit a unified description in terms of category-theoretic
fibrations.

\paragraph{Open problems.}
\begin{itemize}
\item \textbf{Universal property.}  Is the WIS construction the free
Hilbert space generated by the incidence relation?
\item \textbf{Monoidal structure.}  Do WIS admit a tensor product
corresponding to the composition of encryption systems?
\end{itemize}

\section*{Geometric Problems}

Problems concerning the geometry of the posterior manifold, the
UOD, and its relation to existing geometric frameworks.

\section{Curvature and Non-Flat Geometry}\label{ch:sec:open-curvature}

The pullback metric on the posterior manifold is flat because it is
pulled back from Euclidean space.  Replacing the Euclidean target with
a curved space would yield a non-flat geometry whose curvature encodes
additional information about the encoding system.

\paragraph{Open problems.}
\begin{itemize}
\item \textbf{Ricci curvature.}  Does non-negative Ricci curvature of
the posterior manifold imply stronger Lipschitz bounds?
\item \textbf{Comparison geometry.}  How does the UOD behave under
Ricci curvature bounds?
\item \textbf{Spectral geometry.}  Relationship between the spectrum of
the posterior Laplacian and the combinatorial expansion of the
incidence graph?
\end{itemize}

\section{Spectral Theory of the Posterior Laplacian}\label{ch:sec:open-spectral}

The posterior Laplacian $L=A^{\mathsf T}QA$ depends on the
distribution $P$ through the prior weights.  Its spectrum is therefore
a function on the posterior manifold.

\paragraph{Open problems.}
\begin{itemize}
\item \textbf{Eigenvalue bounds.}  What is the range of the smallest
eigenvalue of $L$ as $P$ varies over the simplex?
\item \textbf{Eigenvectors and incognisability.}  Do the eigenvectors
of $L$ correspond to natural partitions of the message space?
\item \textbf{Isoperimetric inequalities.}  Does the spectral gap of
$L$ satisfy a Cheeger-type inequality?
\end{itemize}

\section{Continuous Limits and PDEs}\label{ch:sec:open-continuous-limits}

As message and ciphertext spaces grow large, the discrete WIS
converges to a continuous structure.  The UOD becomes a functional on
probability measures.  Understanding this limit is essential for
applications to location privacy and differential privacy.

\paragraph{Open problems.}
\begin{itemize}
\item \textbf{Well-posedness.}  Does the UOD have a unique minimiser
in the continuous limit?
\item \textbf{Gamma-convergence.}  Do discrete critical points converge
to solutions of the continuous Euler--Lagrange equations?
\item \textbf{Relation to optimal transport.}  Does the UOD gradient
flow converge to a Wasserstein gradient flow?
\end{itemize}

\section{Quantum Extensions}\label{ch:sec:open-quantum}

Generalise WIS to non-commutative probability spaces.  The universal
optimality defect becomes a quantum divergence:
\[
\mathcal D_Q(\rho,\sigma) = \operatorname{Tr}\rho(\log\rho-\log\sigma)^2
- \bigl(\operatorname{Tr}\rho(\log\rho-\log\sigma)\bigr)^2.
\]

\paragraph{Open problems.}
\begin{itemize}
\item \textbf{Quantum Lipschitz bounds.}  Does the quantum UOD satisfy
the same bounds as its classical counterpart?
\item \textbf{Quantum Three-Level Theorem.}  Does the quantum UOD
collapse to finitely many distinct eigenvalues?
\item \textbf{Quantum Apollonian construction.}  Do quantum cell
encoders exist?
\end{itemize}

\section*{Algorithmic Problems}

Problems concerning the computation, optimisation, and numerical
treatment of the UOD and related quantities.

\section{Optimisation of the Universal Optimality Defect}\label{ch:sec:open-optimisation-uod}

The defect $\delta(Y)$ defines a cost function on the space of WIS
weights.  Minimising $\delta$ over admissible weights, subject to
integrality constraints, may yield approximate encoders for systems
where exact universal optimality is unattainable.

\paragraph{Open problems.}
\begin{itemize}
\item \textbf{Algorithmic minimisation.}  Develop efficient algorithms
for minimising the UOD over weight spaces with integrality constraints.
\item \textbf{Trade-off curves.}  Characterise the Pareto frontier
between UOD and key length for families of incidence structures.
\item \textbf{Finite-dimensional reduction.}  Can the three-level
structure be exploited for numerical computation in high dimensions?
\end{itemize}

\section{Gradient Flows and Numerical Optimisation}\label{ch:sec:open-gradient-flows}

The variational formulation of Part~V suggests a natural gradient flow
for the UOD on the posterior manifold:
\[
\frac{dP_t}{dt} = -\nabla_g \mathcal D(P_t, U_n), \qquad
P_0 = P_{\mathrm{init}}.
\]
In WIS coordinates this becomes a system of ODEs whose stationary
points are the critical points of the UOD.  The Three-Level Theorem
implies that any convergent trajectory ends at a configuration with at
most three distinct log-likelihood values.

\paragraph{Open problems.}
\begin{itemize}
\item \textbf{Convergence.}  Does the gradient flow converge for every
initial condition?  The flat metric suggests linear convergence for
convex functionals, but the UOD is not globally convex.
\item \textbf{Bifurcation analysis.}  The pitchfork bifurcations
described in Section~25.9.1 should mark qualitative changes in the
flow landscape.
\item \textbf{Finite-dimensional reduction.}  Can the three-level
structure be exploited for numerical computation in high dimensions?
\end{itemize}

\section*{Applications}

Problems concerning the practical use of the WIS framework in
cryptography, privacy, and information theory.

\paragraph{Privacy--utility trade-offs.}
The UOD is a continuous measure of privacy leakage.  Understanding how
it interacts with utility---the usefulness of the released data---is
essential for practical deployment.  Does every privacy mechanism
admit an optimal trade-off curve parameterised by the UOD?
